\documentclass[11pt, letterpaper]{article}

\usepackage[margin=1in]{geometry}
\usepackage[T1]{fontenc}
\usepackage{lmodern}
\usepackage{parskip}
\usepackage{booktabs}
\usepackage{longtable}
\usepackage{array}
\usepackage{hyperref}
\usepackage{enumitem}
\usepackage{xcolor}
\usepackage{listings}
\usepackage{fancyhdr}

\hypersetup{
    colorlinks=true,
    linkcolor=blue!70!black,
    urlcolor=blue!60!black,
}

\lstset{
    basicstyle=\small\ttfamily,
    breaklines=true,
    frame=single,
    backgroundcolor=\color{gray!8},
    showstringspaces=false,
    tabsize=2,
}

\pagestyle{fancy}
\fancyhf{}
\lhead{Professor Reference Repo Analysis}
\rhead{\texttt{last30days-skill}}
\cfoot{\thepage}

\title{
    \textbf{Professor Reference Repo Analysis} \\[0.4em]
    \Large How \texttt{last30days-skill} Applies to the Polymarket Insider-Signal Project
}
\author{
    Vraj Patel \\
    M.S. Data Science, Stony Brook University \\
    CSE 520 --- Quantitative Finance
}
\date{March 31, 2026}

\begin{document}
\maketitle
\tableofcontents
\newpage

\section{Executive Bottom Line}

The professor repository at \texttt{Prof-refe/last30days-skill} is not a replacement for your Polymarket collector. It solves a different problem: turning a topic into a multi-source research brief by searching several public information sources, normalizing the returned items, ranking them, deduplicating them, saving them, and rendering reusable summaries.

For your project, that means the repo is best understood as an \textbf{intelligence layer template}, not a \textbf{sensing layer template}. Your existing codebase already does the sensing-layer job better because it stores high-frequency Polymarket microstructure. The professor repo is useful because it shows how to:

\begin{enumerate}[nosep]
    \item normalize heterogeneous evidence into one schema,
    \item score and rank evidence from multiple sources,
    \item persist watchlists and repeated runs,
    \item generate compact explainable briefings,
    \item build an evaluation harness around retrieval quality.
\end{enumerate}

\textbf{Best practical use:} do not merge the professor repo wholesale into your collector. Instead, port its architectural patterns into a new evidence-orchestration subsystem that activates \emph{after} your Polymarket detector flags a market, wallet cluster, or suspicious episode.

\section{What the Reference Repo Is Doing}

\subsection{Core Purpose}

\texttt{last30days-skill} is a topic-centric research engine. The user gives it a topic such as a company, person, tool, or event. The system then queries multiple sources from the recent past, scores what it finds, and returns a synthesized research package.

\subsection{Functional Flow}

\begin{lstlisting}[title={High-Level Runtime Flow}]
User topic
  -> query cleanup / query typing
  -> parallel source search
  -> normalization to a common schema
  -> date filtering
  -> scoring and ranking
  -> dedupe and cross-source linking
  -> render compact output
  -> optionally save to watchlist/store/briefing DB
\end{lstlisting}

\subsection{Main Sources It Can Use}

The repo supports a mix of connectors, including Reddit, X, YouTube, TikTok, Instagram, Hacker News, Polymarket, web search, Bluesky, Truth Social, and optional Xiaohongshu. Not every source is always active; the environment/config layer chooses what is available and what fallback path to use.

\subsection{Main Methodological Ideas}

\begin{enumerate}[nosep]
    \item \textbf{Query expansion.} The system does not trust the raw user query alone; it extracts a core subject and expands the query for source search.
    \item \textbf{Canonical normalization.} Source-specific data gets transformed into shared item types with common fields such as title/text, date, engagement, relevance, and score.
    \item \textbf{Multi-factor scoring.} Results are ranked by a weighted mix of relevance, recency, and engagement.
    \item \textbf{Cross-source evidence fusion.} The repo links near-duplicate items and marks when a story appears on multiple sources.
    \item \textbf{Persistence and reruns.} Watchlists, stored findings, and briefings make the tool useful over repeated sessions rather than only as a one-shot search.
    \item \textbf{Evaluation discipline.} Search-quality and synthesis-evaluation scripts show a habit of measuring retrieval quality rather than relying only on subjective impressions.
\end{enumerate}

\section{Why It Is Usable for Your Situation}

\subsection{Where It Fits}

Your current project needs two very different planes:

\begin{enumerate}[nosep]
    \item a \textbf{market-data plane} that captures prices, trades, order books, wallets, and timing, and
    \item an \textbf{evidence plane} that explains what might be happening in the world around a flagged market.
\end{enumerate}

Your current codebase already targets the first plane. The professor repo is relevant because it shows how to implement the second plane cleanly.

\subsection{Why This Matters for Your Hypotheses}

\textbf{H1 (Goodhart's Law).} You need a public-information layer to test whether a market loses predictive power when it becomes highly observed and heavily discussed. The professor repo gives patterns for measuring public attention and discussion intensity.

\textbf{H2 (new-wallet insider signatures).} A wallet-flow anomaly alone is not enough. You also need a context engine that asks: is there already public chatter on X, Google/web, or other channels? The reference repo gives the retrieval, scoring, and briefing logic for that context check.

\textbf{H3 (leading indicator).} To claim that insider-weighted Polymarket signals lead real-world events, you need timestamped external evidence. The professor repo provides the template for collecting, normalizing, and ranking those outside signals.

\subsection{Most Important Conceptual Translation}

The professor repo is \textbf{topic-centric}. Your system must become \textbf{market-centric}, \textbf{wallet-centric}, and \textbf{episode-centric}. That means the reusable parts are mostly \emph{patterns}, not final schemas.

\begin{lstlisting}[title={Mental Model}]
Your collector            = market sensing
Professor repo            = external evidence orchestration
Target thesis system      = anomaly detection + evidence fusion + alerting
\end{lstlisting}

\section{Direct Reuse vs. Adaptation}

\subsection{Directly Reusable Patterns}

\begin{longtable}{p{3.6cm}p{5.0cm}p{5.0cm}}
\toprule
\textbf{Reference File} & \textbf{What It Already Does Well} & \textbf{How To Reuse It} \\
\midrule
\endhead
\texttt{scripts/store.py} & SQLite accumulation, WAL mode, repeated runs, settings, dedup-friendly persistence & Rebuild as the signal/evidence store for alerts, episodes, briefings, analyst actions, and recurring monitored markets \\
\texttt{scripts/watchlist.py} & Repeat-run orchestration and schedule semantics & Adapt into a watchlist for markets, event groups, or wallet clusters instead of free-form topics \\
\texttt{scripts/briefing.py} & Converts stored findings into structured briefing data & Reuse the pattern for analyst briefs, morning summaries, and post-alert recap packets \\
\texttt{scripts/lib/schema.py} & Shared item schema across heterogeneous sources & Reuse the normalization philosophy for X/news/web evidence tied to Polymarket events \\
\texttt{scripts/lib/normalize.py} & Canonicalizes source outputs before downstream use & Reuse as the normalization layer in your evidence subsystem \\
\texttt{scripts/lib/score.py} & Weighted ranking based on relevance/recency/engagement & Reuse the scoring framework, but change the features and weights to event- and signal-aware evidence ranking \\
\texttt{scripts/lib/dedupe.py} & Near-duplicate detection and cross-source linking & Reuse to collapse repeated headlines/posts into one evidence cluster per event \\
\texttt{scripts/lib/render.py} & Human-readable compact summaries with honesty about missing/fresh data & Reuse to generate explainable alerts and analyst-facing summaries \\
\texttt{scripts/evaluate\_search\_quality.py} & Offline retrieval-evaluation harness & Reuse as the model for benchmarking your evidence retrieval and ranking layer \\
\bottomrule
\end{longtable}

\subsection{Reusable Only With Adaptation}

\begin{longtable}{p{3.6cm}p{4.8cm}p{5.2cm}}
\toprule
\textbf{Reference File} & \textbf{Value} & \textbf{Required Adaptation} \\
\midrule
\endhead
\texttt{scripts/lib/polymarket.py} & Good example of free Polymarket search via Gamma plus query expansion & It is event search, not trade/order-flow capture; cannot detect wallets, coordinated entry, or real-time anomalies \\
\texttt{scripts/lib/query.py} & Good heuristics for stripping noise from natural-language prompts & Needs to be replaced by event-title, market-tag, and entity extraction from your own detected episodes \\
\texttt{scripts/lib/query\_type.py} & Useful for source routing based on query class & Must be reworked so routing depends on market category, anomaly type, and urgency rather than topic type alone \\
\texttt{scripts/lib/entity\_extract.py} & Useful for extracting handles, hashtags, and communities from first-pass evidence & Needs event-specific seeding from market title, resolution text, tags, named entities, and related tickers/regions \\
\texttt{scripts/lib/quality\_nudge.py} & Useful concept for measuring evidence completeness & Should become a confidence/completeness score for the alert package, not a setup nudge \\
\texttt{scripts/lib/models.py} and source-search backends & Useful for choosing different LLM/search providers & Must be integrated behind your own evidence-agent interface with stronger reliability controls \\
\bottomrule
\end{longtable}

\subsection{Should Not Be Reused for the Core Collector}

\begin{longtable}{p{3.6cm}p{10.2cm}}
\toprule
\textbf{Reference Surface} & \textbf{Why It Should Not Be Reused Directly} \\
\midrule
\endhead
\texttt{scripts/last30days.py} as a whole & The orchestrator is designed around ad hoc topic research, not always-on market surveillance or market microstructure ingestion \\
\texttt{scripts/lib/polymarket.py} as a collector substitute & It only searches Polymarket events and markets; it does not capture order books, trades, wallet identity, or continuous snapshots \\
\texttt{.claude-plugin/}, \texttt{agents/openai.yaml}, \texttt{gemini-extension.json} & Useful for packaging the skill, but not part of the detection logic itself \\
\texttt{vendor/} and packaged Bird code & Third-party vendored dependency surface; useful operationally, but not something to treat as your research architecture \\
\texttt{assets/}, launch copy, and marketing docs & Productization material rather than a modeling or data-engineering contribution \\
\bottomrule
\end{longtable}

\section{Concrete Architecture You Can Build From It}

\subsection{Recommended Combined Architecture}

\begin{lstlisting}[title={Recommended Combined System}]
Polymarket collector
  -> market snapshots / order books / trades / wallets
  -> anomaly detector
  -> episode builder
  -> evidence orchestrator
       -> X search
       -> web/news search
       -> optional Reddit / YouTube / HN
       -> normalize
       -> score
       -> dedupe
       -> evidence summary
  -> alert ranker
  -> Telegram / Discord / SMS
  -> archive / backtest / analyst feedback store
\end{lstlisting}

\subsection{Best Immediate Ports}

\begin{enumerate}[nosep]
    \item Create a new local package such as \texttt{evidence/} or \texttt{context\_engine/}.
    \item Port the \texttt{schema.py}, \texttt{normalize.py}, \texttt{score.py}, \texttt{dedupe.py}, and \texttt{render.py} ideas into that package.
    \item Replace the professor repo's ``topic'' unit with your own unit: \texttt{signal\_episode}.
    \item Use your signal episode to seed source-specific searches: event title, region, entity names, market tags, and resolution text.
    \item Save all external findings in a dedicated evidence store keyed by \texttt{signal\_id}, \texttt{event\_id}, and retrieval time.
    \item Generate briefings only after the market detector has raised a candidate anomaly.
\end{enumerate}

\subsection{Best Data-Model Translation}

\begin{longtable}{p{4.3cm}p{4.2cm}p{4.8cm}}
\toprule
\textbf{Professor Repo Concept} & \textbf{Current Meaning There} & \textbf{Meaning in Your Project} \\
\midrule
\endhead
\texttt{topic} & User-supplied research subject & Market, event, event-group, or anomaly episode \\
\texttt{research\_run} & One source-collection pass for a topic & One evidence retrieval pass for a flagged signal \\
\texttt{finding} & Retrieved post/article/video/thread & External evidence item linked to a signal episode \\
\texttt{briefing} & Daily or weekly topic digest & Analyst alert packet or periodic portfolio risk summary \\
\texttt{score} & Relevance/recency/engagement rank & Evidence usefulness for validating or invalidating the anomaly \\
\bottomrule
\end{longtable}

\section{Sequential File-by-File Review}

\subsection{Top-Level Authored Files}

\begin{longtable}{p{4.0cm}p{4.9cm}p{4.7cm}}
\toprule
\textbf{File} & \textbf{Role in Repo} & \textbf{Usefulness for Your Project} \\
\midrule
\endhead
\texttt{README.md} & Product-level overview, supported sources, usage framing & High; good summary of intended architecture and source mix \\
\texttt{CHANGELOG.md} & Evolution of features and source additions & Medium; reveals maturity and design priorities \\
\texttt{CLAUDE.md} & Repo structure notes and implementation guidance & Medium; helps navigate architecture quickly \\
\texttt{SPEC.md} & Formal module-level design intent & High; clearest architecture map in the repo \\
\texttt{TASKS.md} & Build checklist and implementation milestones & Medium; useful for seeing what was considered essential \\
\texttt{SKILL.md} & Main assistant-facing runtime instructions & Medium; helpful for routing and invocation design \\
\texttt{SKILL-original.md} & Earlier version of the skill contract & Low; mainly historical context \\
\texttt{release-notes.md} & Product-level summary of shipped capabilities & Low for code reuse; useful for understanding positioning \\
\texttt{.claude-plugin/plugin.json} & Claude packaging manifest & Low for research logic \\
\texttt{.claude-plugin/marketplace.json} & Claude marketplace metadata & Low for research logic \\
\texttt{agents/openai.yaml} & OpenAI/agent interface metadata & Low for core system; useful only if you later package an assistant tool \\
\texttt{gemini-extension.json} & Gemini packaging and env settings & Low for core system \\
\texttt{hooks/scripts/check-config.sh} & Startup diagnostic and configuration banner & Medium; useful pattern for environment health checks \\
\bottomrule
\end{longtable}

\subsection{Primary Scripts}

\begin{longtable}{p{4.0cm}p{4.9cm}p{4.7cm}}
\toprule
\textbf{File} & \textbf{Role in Repo} & \textbf{Usefulness for Your Project} \\
\midrule
\endhead
\texttt{scripts/last30days.py} & Main topic-research orchestrator, connector coordination, fallbacks & Medium as inspiration, low as direct code reuse \\
\texttt{scripts/store.py} & Persistent research/watchlist database with FTS and migrations & Very high; strong pattern to adapt \\
\texttt{scripts/watchlist.py} & Watchlist CLI and recurring rerun entrypoint & High; strong pattern to adapt \\
\texttt{scripts/briefing.py} & Structured briefing-data generator & High; strong pattern to adapt \\
\texttt{scripts/evaluate\_search\_quality.py} & Retrieval evaluation harness & High; directly useful as a benchmark pattern \\
\texttt{scripts/evaluate-synthesis.py} & Manual synthesis-comparison framework & Medium; useful if you later compare alert narratives or analyst briefs \\
\texttt{scripts/generate-synthesis-inputs.py} & Converts stored JSON to compact markdown for synthesis comparison & Medium; useful for your later prompt-evaluation workflow \\
\texttt{scripts/sync.sh} & Installs/syncs the skill into assistant environments & Low for thesis system logic \\
\texttt{scripts/test-v1-vs-v2.sh} & Comparison helper for older skill versions & Low; process artifact \\
\bottomrule
\end{longtable}

\subsection{Library Modules}

\begin{longtable}{p{4.0cm}p{4.6cm}p{5.0cm}}
\toprule
\textbf{File} & \textbf{Primary Role} & \textbf{Reuse Verdict} \\
\midrule
\endhead
\texttt{scripts/lib/env.py} & Environment loading, key discovery, cookie-domain registry & Reuse pattern only \\
\texttt{scripts/lib/cache.py} & Query-result and model cache helpers & Reuse pattern only \\
\texttt{scripts/lib/http.py} & Retry-aware HTTP utilities & Reuse pattern only \\
\texttt{scripts/lib/models.py} & Model/provider selection helpers & Reuse with adaptation \\
\texttt{scripts/lib/query.py} & Query cleanup and core-subject extraction & Reuse with adaptation \\
\texttt{scripts/lib/query\_type.py} & Query classification and source prioritization & Reuse with adaptation \\
\texttt{scripts/lib/schema.py} & Canonical typed schema for normalized items & Strong reuse candidate \\
\texttt{scripts/lib/dates.py} & Date parsing and recency helpers & Strong reuse candidate \\
\texttt{scripts/lib/relevance.py} & Text-overlap relevance heuristics & Reuse with adaptation \\
\texttt{scripts/lib/normalize.py} & Converts source-specific payloads to canonical types & Strong reuse candidate \\
\texttt{scripts/lib/score.py} & Computes weighted ranking subscores and final scores & Strong reuse candidate \\
\texttt{scripts/lib/dedupe.py} & Near-duplicate detection and cross-source linking & Strong reuse candidate \\
\texttt{scripts/lib/render.py} & Compact human-readable rendering of normalized reports & Strong reuse candidate \\
\texttt{scripts/lib/polymarket.py} & Gamma search over Polymarket events and markets & Limited use; not a collector substitute \\
\texttt{scripts/lib/websearch.py} & Native web search parsing and normalization & Strong evidence-layer pattern \\
\texttt{scripts/lib/parallel\_search.py} & Parallel AI web search backend wrapper & Optional provider adapter \\
\texttt{scripts/lib/brave\_search.py} & Brave search backend adapter & Optional provider adapter \\
\texttt{scripts/lib/exa\_search.py} & Exa search backend adapter & Optional provider adapter \\
\texttt{scripts/lib/openrouter\_search.py} & OpenRouter/Perplexity web backend adapter & Optional provider adapter \\
\texttt{scripts/lib/openai\_reddit.py} & Reddit discovery through OpenAI responses/web search & Optional source connector \\
\texttt{scripts/lib/reddit\_public.py} & Free Reddit JSON search fallback & Optional source connector \\
\texttt{scripts/lib/reddit.py} & ScrapeCreators-backed Reddit search & Optional source connector \\
\texttt{scripts/lib/reddit\_enrich.py} & Fetches thread comments and extracts discussion insights & Strong pattern for evidence enrichment \\
\texttt{scripts/lib/xai\_x.py} & X search via xAI & Strong evidence-layer pattern \\
\texttt{scripts/lib/bird\_x.py} & Free local X search via bundled Bird client & Strong evidence-layer pattern \\
\texttt{scripts/lib/scrapecreators\_x.py} & X search via ScrapeCreators API & Optional evidence connector \\
\texttt{scripts/lib/youtube\_yt.py} & YouTube search and transcript extraction & Optional evidence connector \\
\texttt{scripts/lib/tiktok.py} & TikTok search/normalize helpers & Optional, lower priority \\
\texttt{scripts/lib/instagram.py} & Instagram search/normalize helpers & Optional, lower priority \\
\texttt{scripts/lib/bluesky.py} & Bluesky connector & Optional, lower priority \\
\texttt{scripts/lib/truthsocial.py} & Truth Social connector & Optional, lower priority \\
\texttt{scripts/lib/hackernews.py} & Hacker News retrieval and scoring & Optional but interesting for tech-linked events \\
\texttt{scripts/lib/entity\_extract.py} & Extracts handles, hashtags, subreddits from first-pass hits & Useful adaptation target \\
\texttt{scripts/lib/quality\_nudge.py} & Coverage/completeness scoring and user nudge generation & Useful adaptation target \\
\texttt{scripts/lib/setup\_wizard.py} & First-run setup helper & Low for thesis logic \\
\texttt{scripts/lib/ui.py} & Terminal spinner/banner utilities & Low for thesis logic \\
\texttt{scripts/lib/chrome\_cookies.py} & Chrome cookie extraction for auth bootstrap & Operational helper only \\
\texttt{scripts/lib/cookie\_extract.py} & Cross-browser cookie extraction wrapper & Operational helper only \\
\texttt{scripts/lib/safari\_cookies.py} & Safari cookie extraction & Operational helper only \\
\texttt{scripts/lib/xiaohongshu\_api.py} & Xiaohongshu REST connector & Optional, low priority \\
\bottomrule
\end{longtable}

\subsection{Open Variant and Reference Docs}

\begin{longtable}{p{4.0cm}p{4.9cm}p{4.7cm}}
\toprule
\textbf{File} & \textbf{Role in Repo} & \textbf{Usefulness for Your Project} \\
\midrule
\endhead
\texttt{variants/open/SKILL.md} & Multi-mode routing for research, watchlist, briefing, and history & Medium; useful for thinking about operator workflows \\
\texttt{variants/open/context.md} & Lightweight persistent memory for preferences and source notes & Medium; could inspire analyst preference or strategy memory \\
\texttt{variants/open/references/research.md} & Instructions for one-shot research mode & Medium; workflow reference only \\
\texttt{variants/open/references/watchlist.md} & Watchlist behavior contract & Medium; workflow reference only \\
\texttt{variants/open/references/briefing.md} & Briefing behavior contract & Medium; workflow reference only \\
\texttt{variants/open/references/history.md} & History-query behavior contract & Medium; workflow reference only \\
\texttt{docs/how-search-works.md} & Detailed explanation of Reddit/X search architecture & High; useful for understanding the repo's reasoning model \\
\texttt{docs/search-quality-eval.md} & Describes retrieval-quality evaluation methodology & High; useful for adapting evaluation practices \\
\texttt{docs/pr-credits.md} & Contribution/process notes & Low \\
\texttt{docs/v2.1-launch-copy.md}, \texttt{docs/v2.1-tweets.md}, \texttt{docs/v2.5-launch-tweets.md} & Launch/marketing materials & Low \\
\bottomrule
\end{longtable}

\subsection{Tests, Fixtures, and Supporting Surfaces}

\begin{longtable}{p{4.0cm}p{9.8cm}}
\toprule
\textbf{Surface} & \textbf{What It Tells You} \\
\midrule
\endhead
\texttt{tests/test\_schema\_roundtrip.py} and \texttt{tests/test\_normalize.py} & The repo takes canonical schema stability seriously; this is exactly the discipline your evidence subsystem should adopt \\
\texttt{tests/test\_score.py}, \texttt{tests/test\_relevance.py}, \texttt{tests/test\_dedupe.py} & Ranking and deduplication are treated as first-class logic rather than ad hoc formatting \\
\texttt{tests/test\_query.py}, \texttt{tests/test\_query\_type.py}, \texttt{tests/test\_entity\_extract.py} & Query formation and source routing are explicitly tested, which is important if you later build event-aware evidence search \\
\texttt{tests/test\_evaluate\_search\_quality.py} & The evaluation harness itself is tested, which is a strong signal of engineering maturity \\
\texttt{tests/test\_polymarket.py} & Confirms that the repo treats Polymarket as a search/discovery source rather than a market microstructure source \\
\texttt{tests/test\_bird\_x.py}, \texttt{tests/test\_scrapecreators\_x.py}, \texttt{tests/test\_youtube\_yt.py}, \texttt{tests/test\_reddit\_public.py}, \texttt{tests/test\_reddit\_enrich.py}, and others & Connector-level tests show the external-source layer is modular and reasonably mature \\
\texttt{fixtures/} & Sample payloads used for deterministic testing; useful pattern for your own evidence and alert fixtures \\
\texttt{docs/plans/} and \texttt{docs/test-results/} & Historical planning and evaluation artifacts; good for repo context, but not runtime code to reuse \\
\texttt{vendor/} & Third-party packaged Bird code and related assets; operational dependency, not your target design pattern \\
\bottomrule
\end{longtable}

\section{What You Should Actually Do With It}

\subsection{Recommended Implementation Strategy}

\begin{enumerate}[nosep]
    \item Keep your current Polymarket ingestion repo as the authoritative source for markets, trades, order books, and wallet-derived episodes.
    \item Create a new evidence package inside your repo inspired by the professor repo's normalization, scoring, dedupe, storage, and briefing layers.
    \item Use that package only after a detector identifies a candidate event or suspicious coordinated wallet pattern.
    \item Store external evidence point-in-time so that backtests can replay what was known when the signal fired.
    \item Borrow the professor repo's evaluation mindset and create a benchmark for evidence retrieval quality, alert relevance, and briefing usefulness.
\end{enumerate}

\subsection{Specific Modules Worth Re-Implementing First}

\begin{enumerate}[nosep]
    \item \textbf{Evidence schema} inspired by \texttt{schema.py}
    \item \textbf{Evidence normalization} inspired by \texttt{normalize.py}
    \item \textbf{Evidence ranking} inspired by \texttt{score.py}
    \item \textbf{Evidence dedupe and cross-linking} inspired by \texttt{dedupe.py}
    \item \textbf{Signal/evidence store} inspired by \texttt{store.py}
    \item \textbf{Alert packet rendering} inspired by \texttt{render.py} and \texttt{briefing.py}
    \item \textbf{Retrieval evaluation harness} inspired by \texttt{evaluate\_search\_quality.py}
\end{enumerate}

\subsection{What Not To Waste Time Porting Early}

\begin{enumerate}[nosep]
    \item assistant packaging files,
    \item setup-wizard UX,
    \item social sources beyond X and web/news,
    \item the repo's own Polymarket search connector as a substitute for real market ingestion,
    \item marketing and release-note materials.
\end{enumerate}

\section{Final Assessment}

The professor repo is highly usable for your situation, but only if used in the right place. It should \textbf{not} replace your current Polymarket pipeline. It \textbf{should} shape the subsystem that turns a suspicious market episode into an explainable research object.

That is the key takeaway:

\begin{enumerate}[nosep]
    \item your repo should own \textbf{market sensing},
    \item the professor repo should inspire \textbf{context fusion},
    \item your thesis system should combine both into \textbf{ranked real-time alerts}.
\end{enumerate}

\end{document}
